\subsubsection{A nonsplit extension of the two blocks}
\label{par:full-block-extension}
\hypertarget{source-cohomology-full-block-extension}{}

To apply the two-link test to a point with a full block, one first replaces the split representation by a nonsplit extension. The replacement must preserve the characteristic polynomials on which that test depends.

\begin{proposition}[Full-block extension with the same characteristic polynomials]\label{prop:full-block-extension}
For the normalized two-block representation with one full block
and the marked ordinary flag of the rank-four calculation, write \(r\) for the normalized representation. The small complete integral model admits a nonsplit extension \(r'\) with the same similitude and ordinary diagonal characters and with 
\[
 \det(X-r'(g))=\det(X-r(g))\qquad\text{for every }g.
\]
 It satisfies the same prescribed inertia conditions
and can be used in the nonsplit branch in either characteristic.
\end{proposition}

\begin{proof}
\proofstep{Projection preserving the flag and character polynomials.}

The alternating form pairs the two planes
\[
          V_-=\langle e_0,e_2\rangle,\qquad
          V_+=\langle e_1,e_3\rangle .
\]
On each block complex conjugation is the scalar with its residual
sign. Indeed its square is one and subtracting the opposite scalar
gives an invertible matrix, since \(\ell\) is odd.
In the one-sided triangular case, first take the block-diagonal
projection. This is multiplicative on the block-triangular algebra
and commutes with the symplectic adjoint. It preserves the multiplier,
the two block matrices and the total characteristic polynomial. In
the normalized chart the chosen local element \(\phi\) is block
diagonal and its two residual block polynomials are coprime.

If \(f_-f_+\) is the product of the two monic factors
at \(r(\phi)\), choose \(af_-+bf_+=1\) by coprimality.
Then
\[
                      e_-=b(r(\phi))f_+(r(\phi))
\]
projects onto \(V_-\).
It preserves the ordinary flag. Write \(a\in\{0,1\}\) for the action of \(e_-\) on a flag quotient. If \(X\) acts there by \(\lambda\), the projected matrix acts by
\[
 a^2\lambda+(1-a)^2\lambda=\lambda .
\]
Thus the ordinary diagonal characters are preserved. Every prescribed inertia element that acted
as the identity still acts as the identity. We use this projected
point in characteristic \(\ell\), and the given block-diagonal point
in the characteristic-zero full-block test.

\proofstep{The rank-three off-block module.}

Use the small integral model of this point, with adically
continuous matrices and flag over a complete DVR \(A\), of finite
residue field of characteristic \(\ell\), and a local injection
into the ambient valuation ring. In particular it is this model
that is complete at a uniformizer \(\pi\), in either characteristic.
Write \(r\) for its block-diagonal representation, \(u\) for
the symplectic flag basis, and \(s\) for the selected element of
the rational decomposition group \(G_\ell\) with
\(\epsilon_\ell(s)=1\). Its residual eigenvalues separate the two
blocks. Define
\[
 \begin{aligned}
 M&=\{X: X(V_+)=0,\ X(V_-)\subset V_+,\ X^*=-X\},\\
 P&=\{X\in M:u^{-1}Xu\text{ is strictly upper triangular}\}.
 \end{aligned}
\]
Here \(*\) is the symplectic adjoint and \(G_\Q\) acts on \(M\) by
\(\rho(g)X=r(g)Xr(g)^{-1}\). The paired-block description identifies
\(M\) with the symmetric bilinear forms on \(V_-\), so it is free
of rank three.

All products of two elements of \(M\) vanish.
Complex conjugation \(c\) acts as minus one, and \(\rho(s)-1\) is
invertible: on matrices from \(V_-\) to \(V_+\) its residual
eigenvalues are ratios of residual eigenvalues from distinct
blocks minus one. Since \(P\) is saturated, reduction modulo \(\pi\) preserves
the short exact sequence with quotient \(M/P\). For \(T=\rho(s)-1\), a basis adapted to \(P\subset M\) gives
\[
 \det(T|_M)=\det(T|_P)\det(T|_{M/P})\in A^\times .
\]
Over the local ring \(A\) both factors are units, so both induced maps are invertible. Stability under \(G_\ell\) follows because upper
triangular changes of basis preserve the strictly upper triangular
matrices. Saturation follows by testing the lower and diagonal
entries of \(u^{-1}Xu\) over the domain \(A\).

\proofstep{A nonzero ordinary local submodule.}

The submodule \(P\) is nonzero. The residual labels of the first
two flag quotients are in \(V_-\) and \(V_+\), respectively. The
block projectors of \(r(s)\) preserve the flag by the same
coprime-factor argument. Its Lagrangian plane \(F_2\) therefore
contains a rank-one summand in \(V_+\). For a nonzero vector
\(w\) in that summand, the endomorphism
\(x\mapsto\langle w,x\rangle w\) belongs to \(M\). Isotropy gives \(X(F_2)=0\), while the image lies in \(Aw\subset F_2\). Thus
\[
 X(F_1)=X(F_2)=0,\qquad X(F_3)\subset F_2,\qquad X(F_4)\subset F_2\subset F_3 .
\]
This is the required strict flag condition.

\proofstep{Produce a normalized ordinary cocycle.}

A normalized nonsplit extension requires a nonzero local class with the prescribed ordinary behavior. This is the role of the local submodule in the global cocycle construction.


To obtain the extension, put \(M_n=M/\pi^n M\) and let
\(P_n\subset M_n\) be the image of \(P\). Use the Greenberg--Wiles
cardinality formula for finite \(\ell\)-primary modules of
\(G_{\Q,S}\), where \(S\) contains \(\ell\) and the ramified primes
of \(r\), with local Tate-orthogonal conditions on the dual
module. At a finite prime \(v\nmid\ell\) put \(B=M_n^{I_v}\) and \(T=\operatorname{Frob}_v-1\). For the unramified local condition, finiteness of \(B\) gives
\[
 |H^1_{\rm unr}(G_v,M_n)|=|B/TB|=|\ker(T|_B)|=|H^0(G_v,M_n)|.
\] At \(\ell\) use the image of \(H^1(G_\ell,P_n)\).
The relevant sequence has
\[
     H^0(G_\ell,M_n/P_n)=0\longrightarrow H^1(G_\ell,P_n)
                        \longrightarrow H^1(G_\ell,M_n).
\]
Therefore the chosen local space is isomorphic to
\(H^1(G_\ell,P_n)\). Both \(M_n\) and its Tate dual have
trivial invariants globally and at \(\ell\): test at \(s\), where
the cyclotomic twist is trivial and \(\rho(s)-1\) is invertible.

The same test gives trivial invariants for \(P_n\) and its Tate
dual. The local Euler characteristic and local Tate duality give
\(\lvert H^1(G_\ell,P_n)\rvert=\lvert P_n\rvert\). Write
\(L_v\) for the local conditions and \(L^\perp\) for their
Tate orthogonals on \(M_n^\vee(1)\), the Tate dual. The
cardinality identity is
\[
 \frac{|H^1_L(G_{\Q,S},M_n)|}
 {|H^1_{L^\perp}(G_{\Q,S},M_n^\vee(1))|}
 =\frac{|H^0(G_{\Q,S},M_n)|}{|H^0(G_{\Q,S},M_n^\vee(1))|}
 \prod_{v\in S\cup\{\infty\}}\frac{|L_v|}{|H^0(G_v,M_n)|}
 =|P_n|.
\]
The real place contributes a factor one since \(c\) acts
by \(-1\) and \(\ell\) is odd. In particular
\(|H^1_L(G_{\Q,S},M_n)|\geq |P_n|\). The cardinality formula
and finiteness follow from the finite-module Poitou--Tate sequence,
local duality and the global Euler characteristic.
\cite[Chapter I, Theorems 4.10 and 5.1]{milne-duality}

Every such class has exactly one representative with \(z_n(c)=0\),
by adding a coboundary using the invertibility of two. It vanishes on
\(I_p\cap\ker r\) for \(p\ne\ell\): its restriction to inertia is
a coboundary, and \(\ker r\) acts trivially on coefficients. Its
local condition is
\[
 z_n(h)-(\rho(h)y_n-y_n)\in P_n\qquad(h\in G_\ell)
\]
for some \(y_n\in M_n\). Consider all continuous cocycles with these
three conditions. They form a finite group for each \(n\), since
they are normalized and unramified outside \(S\), and its size is
at least \(\lvert P_n\rvert\). If at some level every such cocycle
vanished modulo \(\pi\), this would also hold at every higher
level by reduction. A cocycle vanishing modulo \(\pi\) at level
\(m\) can be divided by \(\pi\) to give a cocycle at level \(m-1\)
with the same conditions.

At \(s\) reduce the local relation modulo \(\pi\) and project
onto \(M_1/P_1\). When the cocycle reduces to zero, it gives
\[
       (\rho(s)-1)[y_m]=0\text{ in }M_1/P_1 .
\]
Invertibility forces \(y_m\bmod\pi\in P_1\).
Subtract a lift in \(P_m\) from \(y_m\) and divide by \(\pi\).
Saturation of \(P\) then permits division of the whole local
condition. If \(\mathcal Z_m\) denotes the finite group of normalized
cocycles at depth \(m\), then on its subgroup of zero
reduction the division operation just constructed is an
injection into \(\mathcal Z_{m-1}\), with recovery by
multiplication by \(\pi\).

If every cocycle at some depth
\(n\) reduced to zero, this would hold at every \(m\ge n\).
The injections at successive depths would give, for every \(m\ge n\),
\[
 |P_m|\le|\mathcal Z_m|\le|\mathcal Z_{m-1}|\le\cdots\le
 |\mathcal Z_{n-1}|,
\]
with the zero module at depth zero. This contradicts \(|P_m|\to\infty\).
Thus there are cocycles with nonzero reduction at every level.
Take the finite sets of pairs \((z_n,y_n)\) with this property.
Reduction maps preserve the property. The inverse limit is nonempty
by compactness, since any finite set of compatibility conditions
can be met by reducing from a sufficiently high level. Completeness
gives a nonzero continuous cocycle \(z\) over \(A\) and a local vector \(y\) satisfying
\(z(h)-(\rho(h)y-y)\in P\) for all \(h\in G_\ell\). Multiply both by \(\pi\) to make the correction
zero modulo \(\pi\) while preserving nonvanishing over \(A\).

Now set \(r'(g)=(1+z(g))r(g)\) and \(u'=(1-y)u\).
Because products of two elements of \(M\) vanish,
\[
 r'(g)r'(h)
  =(1+z(g))(1+\rho(g)z(h))r(gh)
  =(1+z(g)+\rho(g)z(h))r(gh)
  =r'(gh).
\]
Also \((1+X)^*(1+X)=1+(X+X^*)+X^*X=1\) for \(X\in M\), so
the multiplier stays fixed and \(u'\) is symplectic. At \(h\in
G_\ell\) we have
\[
 (u')^{-1}r'(h)u'=
 \bigl(1+u^{-1}(z(h)-\rho(h)y+y)u\bigr)\,u^{-1}r(h)u.
\]
This is upper triangular with the original diagonal. The inertia
conditions away from \(\ell\) hold because \(z\) kills the prescribed
subgroups of \(\ker r\). The residual representation, local ordinary
characters and small integral model satisfy all the original
point conditions. The two original diagonal blocks are unchanged,
so \(r'\) and \(r\) have identical characteristic polynomials
on \(G_\Q\).

A stable complement to \(V_+\) over the ambient field would be
the graph of a map \(Y:V_-\to V_+\). Write \(r_-(g)\) and
\(r_+(g)\) for the original block matrices. Invariance of
the graph means
\[
        Y r_-(g)=r_+(g)Y+z(g)r_-(g).
\]
At complex conjugation \(z(c)=0\) and the two signs are opposite,
so \(2Y=0\). The residue characteristic is odd, hence \(Y=0\).
The displayed equality for every \(g\) then gives \(z=0\),
contrary to the construction.
Fullness of the minus block implies fullness of the plus block:
the latter is its dual times the multiplier, which acts by
nonzero scalars, and inverse matrices run over the same group.
We may therefore apply the nonsplit block-extension branch of the
modularity argument to \(r'\). The equality of characteristic
polynomials transports its special-fiber or Frobenius conclusions
back to the original block point. \label{cor:cohomology-full-block-extension-1}
\end{proof}
